\section{Functors, naturality, and the meaning of schematic categories}
\label{sec:functors-yoneda-foundation}

Throughout this section, $C$ and $D$ are fixed arithmetic presentations of
categories.  Objects and arrows have natural-number codes.  Each presentation
has arithmetic formulas for validity and arrow equality, and fixed primitive
recursive operations for identities and composition.  We write $f:a\to b$
for the full validity condition and $f\sim g$ for its hom-setoid equality.
The category laws include preservation of typing and of $\sim$ by every
operation.  In particular, $\sim$ need not be numerical equality or a decidable
relation.  All equations below mean this specified equality, and all variables
are accompanied by their displayed typing hypotheses.  Operations may take
the displayed endpoint codes as additional arguments; those arguments are
suppressed in categorical notation.

\subsection{Functors and transformations}

\begin{definition}[Presented functor]
A functor $F:C\to D$ consists of fixed primitive recursive functions $F_0$
and $F_1$ on object and arrow codes.  Its laws say that $F_0$ preserves object
validity, that $f:a\to b$ implies $F_1(f):F_0(a)\to F_0(b)$, that $F_1$
respects arrow equality, and that
\[
 F_1(\id_a)\sim\id_{F_0(a)},\qquad
 F_1(g\circ f)\sim F_1(g)\circ F_1(f).
\]
We suppress the subscripts.  The laws are arithmetic formulas, proved in
$\ipra$ for a concrete presentation, or explicit premises of a relative
structural theorem.
\end{definition}

\begin{definition}[Presented natural transformation]
For fixed functors $F,G:C\to D$, a natural transformation $\alpha:F\Rightarrow G$
is a fixed primitive recursive component function $a\mapsto\alpha_a$ such that
$\alpha_a:Fa\to Ga$ and, for every $f:a\to b$,
\begin{equation}
 G(f)\circ\alpha_a\sim\alpha_b\circ F(f).
 \label{eq:ordinary-naturality}
\end{equation}
Two transformations are equal when their components are equal at every valid
object.  Thus transformation equality is itself an arithmetic universal
formula; it is not a decision procedure for functions.
\end{definition}

\begin{proposition}[The elementary naturality calculus]
Identity functors, composites of functors, identity transformations, vertical
composition, and whiskering are given by primitive recursive operations on
their displayed data.  They satisfy the category and interchange laws
schematically.
\end{proposition}
\begin{proof}
The identity functor leaves each code unchanged.  A composite functor applies
its two fixed object and arrow functions successively.  Preservation of
identity and composition follows by applying the outer functor's congruence
law to the inner functor's corresponding law, and then applying the outer
law.  Associativity of these composites is equality of the displayed terms.

The identity transformation on $F$ has component $\id_{Fa}$.  Its naturality
square commutes by the two identity laws.  If $\alpha:F\Rightarrow G$ and
$\beta:G\Rightarrow H$, define
$(\beta\alpha)_a=\beta_a\circ\alpha_a$.  Its naturality is the calculation
\[
\begin{split}
 H(f)\circ(\beta_a\circ\alpha_a)
 &\sim (H(f)\circ\beta_a)\circ\alpha_a\\
 &\sim (\beta_b\circ G(f))\circ\alpha_a\\
 &\sim\beta_b\circ(G(f)\circ\alpha_a)\\
 &\sim\beta_b\circ(\alpha_b\circ F(f))\\
 &\sim(\beta_b\circ\alpha_b)\circ F(f).
\end{split}
\]
The identity and associativity laws for vertical composition hold component
by component.  Postwhiskering by $K:D\to E$ has components $K(\alpha_a)$;
its laws follow from functoriality.  Prewhiskering by $L:B\to C$ has components
$\alpha_{L b}$; its laws are instances of naturality for arrows $L(f)$.

For clarity, horizontal composition is also explicit.  If
$\alpha:F\Rightarrow G:C\to D$ and $\beta:H\Rightarrow K:D\to E$, put
\[
 (\beta*\alpha)_a=\beta_{Ga}\circ H(\alpha_a).
\]
Naturality of $\beta$ at $\alpha_a$ identifies this with
$K(\alpha_a)\circ\beta_{Fa}$.  Suppose also
$\alpha':G\Rightarrow J$ and $\beta':K\Rightarrow L$.  The component of the
vertical composite of the two horizontal composites is
\[
 \beta'_{Ja}\circ K(\alpha'_a)\circ\beta_{Ga}\circ H(\alpha_a).
\]
Naturality of $\beta$ at $\alpha'_a$ changes the middle two factors to
$\beta_{Ja}\circ H(\alpha'_a)$.  Functoriality of $H$ then gives
\[
 (\beta'_{Ja}\circ\beta_{Ja})\circ H(\alpha'_a\circ\alpha_a),
\]
which is the component of $(\beta'\beta)*(\alpha'\alpha)$.  Parentheses may
be changed by associativity.  This proves interchange.  Every argument just
given is a finite equational derivation, with no use of induction or excluded
middle.
\end{proof}

This result does \emph{not} assert the existence inside arithmetic of a
category whose arrows are all primitive recursive functors or all primitive
recursive transformations.  It supplies a theorem for each fixed finite
configuration of such data.  A separately specified, uniformly evaluable
coded family can itself be made into a presented category if the required
closure and evaluation operations are supplied.  The word ``schematic''
will retain this meaning below.

\subsection{Setoid-valued presheaves without a universe of sets}

\begin{definition}[Presented presheaf]
A presheaf $P$ on $C$ consists of the following data:
\begin{enumerate}
\item an arithmetic formula $P_0(a,x)$ specifying valid elements in its fibre
      at $a$;
\item an arithmetic equivalence relation $x\approx_a y$ on that fibre;
\item a fixed primitive recursive action function, written $P(f)x$, which
      takes $f:a\to b$ and $x\in P(b)$ to an element of $P(a)$.
\end{enumerate}
The action respects both arrow equality and fibre equality, and satisfies
\begin{equation}
 P(\id_a)x\approx_a x,\qquad
 P(g\circ f)x\approx_a P(f)(P(g)x)
 \quad(f:a\to b, g:b\to c, x\in P(c)).
 \label{eq:presheaf-laws}
\end{equation}
A transformation $\alpha:P\Rightarrow Q$ is a fixed primitive recursive
function $(a,x)\mapsto\alpha_a(x)$ which preserves membership and fibre
equality and satisfies
\begin{equation}
 Q(f)(\alpha_b(x))\approx_a\alpha_a(P(f)x).
 \label{eq:presheaf-naturality}
\end{equation}
Transformations are identified pointwise modulo the target fibre relation.
\end{definition}

There is no assumption that $P(a)$ is a separately existing set object.  The
notation means a predicate on natural numbers with a specified equivalence
relation.  Likewise, writing $P:C^{\mathrm{op}}\to\Set$ is at most a
mnemonic for the data above.  The construction never invokes a universe
containing all sets or all maps.

For an object $A$, the representable presheaf $h_A$ is particularly simple:
\[
 h_A(a)=\Hom_C(a,A),\qquad h_A(f)(u)=u\circ f.
\]
Its membership and equality formulas are the hom formulas of $C$.  Its action
is the fixed composition operation of $C$.  Associativity and identity in
$C$ prove its presheaf laws.  The dependence on the parameter $A$ is uniform:
$A$ is an ordinary numerical parameter, not a code for an arbitrary program.

\section{Yoneda, universal elements, and uniqueness of representations}
\label{sec:yoneda}

\begin{theorem}[Primitive recursive schematic Yoneda lemma]
\label{thm:yoneda-pr}
For every presented presheaf $P$, object $A$, and presented transformation
$\alpha:h_A\Rightarrow P$, define
\[
 \operatorname{ev}_A(\alpha)=\alpha_A(\id_A).
\]
Conversely, for $x\in P(A)$ define a transformation
$\widehat{x}:h_A\Rightarrow P$ by
\begin{equation}
 \widehat{x}_a(f)=P(f)x,\qquad f:a\to A.
 \label{eq:yoneda-construction}
\end{equation}
These constructions respect the relevant equivalence relations and satisfy
\[
 \operatorname{ev}_A(\widehat{x})\approx_A x,
 \qquad
 \widehat{\operatorname{ev}_A(\alpha)}_a(f)
   \approx_a\alpha_a(f).
\]
For each fixed input presentation and transformation term, evaluation is a
primitive recursive term, and reconstruction uses the fixed action term with
the extra numerical parameter $x$.  Evaluation is not a universal evaluator
whose input is the code of an arbitrary program.  Thus the customary notation
$\operatorname{Nat}(h_A,P)\cong P(A)$ has an explicit, first-order schematic
meaning in $\ipra$.
\end{theorem}
\begin{proof}
The identity $\id_A$ is a valid element of $h_A(A)$, so preservation of
membership by $\alpha$ makes its evaluation a valid member of $P(A)$.

For the converse, if $f:a\to A$, action typing gives $P(f)x\in P(a)$.
If $f\sim f'$, action congruence gives $P(f)x\approx_a P(f')x$.
For $k:b\to a$, the required naturality equation is exactly
\[
 P(k)(\widehat{x}_a(f))
 =P(k)(P(f)x)
 \approx_b P(f\circ k)x
 =\widehat{x}_b(h_A(k)(f)).
\]
This is the composition law of the presheaf, used in its symmetric direction.
If $x\approx_A y$, action congruence makes $\widehat{x}$ and $\widehat{y}$
pointwise equal.  If $\alpha$ and $\beta$ are pointwise equal, their values
at $(A,\id_A)$ are equal, so evaluation also respects equality.

The first inverse law is
$\operatorname{ev}_A(\widehat{x})=P(\id_A)x\approx_A x$.
For the second, use naturality of $\alpha$ at $f:a\to A$ and apply it to
$\id_A\in h_A(A)$.  It says
\[
 P(f)(\alpha_A(\id_A))
 \approx_a\alpha_a(h_A(f)(\id_A))
 =\alpha_a(\id_A\circ f)
 \approx_a\alpha_a(f).
\]
The final step is the identity law in $C$, followed by congruence of
$\alpha_a$.  This is exactly the required pointwise equality.  No element
has been chosen by search, and no equality was tested.  The witness term
is the action function of $P$ applied to its two displayed arguments.
\end{proof}

\begin{proposition}[Naturality of Yoneda in both variables]
For $\tau:P\Rightarrow Q$, the Yoneda correspondences commute with applying
$\tau$.  For an arrow $u:A\to B$, they commute with precomposition by the
transformation $h_u:h_A\Rightarrow h_B$, where
$(h_u)_a(f)=u\circ f$, and with the action map $P(u):P(B)\to P(A)$.
\end{proposition}
\begin{proof}
The first assertion at evaluation is literal:
\[
 \operatorname{ev}_A(\tau\alpha)
 =\tau_A(\alpha_A(\id_A)).
\]
At reconstruction it is naturality of $\tau$:
\[
 \tau_a(P(f)x)\approx_a Q(f)(\tau_A(x)).
\]
For the second assertion, let $\alpha:h_B\Rightarrow P$.  Evaluation of
$\alpha h_u$ gives
\[
 \alpha_A(u\circ\id_A)\approx_A\alpha_A(u)
 \approx_A P(u)(\alpha_B(\id_B)),
\]
where the last equality is the second inverse identity already proved.
Equivalently, starting with $x\in P(B)$, reconstruction after applying $P(u)$
has component
\[
 P(f)(P(u)x)\approx_a P(u\circ f)x,
\]
which is reconstruction followed by precomposition with $h_u$.
\end{proof}

\begin{theorem}[The Yoneda embedding is schematically fully faithful]
\label{thm:yoneda-fully-faithful}
The association $A\mapsto h_A$, $u\mapsto h_u$ preserves identities and
composition.  Every fixed presented transformation $\alpha:h_A\Rightarrow h_B$
has the unique inducing arrow, modulo hom equality,
\[
 u=\alpha_A(\id_A):A\to B.
\]
In particular, if $h_u$ and $h_v$ are pointwise equal, then $u\sim v$.
\end{theorem}
\begin{proof}
The component of $h_{\id_A}$ at $f$ is $\id_A\circ f\sim f$.
The component of $h_v h_u$ is $v\circ(u\circ f)\sim(v\circ u)\circ f$,
which is the component of $h_{v\circ u}$.  Thus the association is functorial
in the schematic sense just explained.

Apply Theorem~\ref{thm:yoneda-pr} with $P=h_B$.  The reconstruction of
$u\in h_B(A)$ is
$\widehat{u}_a(f)=h_B(f)(u)=u\circ f$, precisely $h_u$.  Its second inverse
identity gives $\alpha_a(f)\sim u\circ f$.  If also
$\alpha_a(f)\sim v\circ f$ for all $a,f$, put $a=A$ and $f=\id_A$.
The identity laws give $u\sim v$.  This proves fullness and faithfulness
without a general function space or any choice operation.
\end{proof}

\begin{corollary}[Isomorphism is detected by representables]
If presented transformations $\alpha:h_A\Rightarrow h_B$ and
$\beta:h_B\Rightarrow h_A$ are pointwise inverse, the arrows
$u=\alpha_A(\id_A)$ and $v=\beta_B(\id_B)$ satisfy
$v\circ u\sim\id_A$ and $u\circ v\sim\id_B$.
\end{corollary}
\begin{proof}
By fullness, $\alpha=h_u$ and $\beta=h_v$ pointwise.  Functoriality turns
$\beta\alpha=\id$ into $h_{v\circ u}=h_{\id_A}$.  Faithfulness gives
$v\circ u\sim\id_A$.  The other identity is obtained from
$\alpha\beta=\id$ in the same way.  Every arrow used in the conclusion has
already been displayed.
\end{proof}

\subsection{Universal elements and the category of elements}

The category of elements $\int P$ has object codes $(a,x)$ with
$a\in C$ and $x\in P(a)$.  A morphism
\[
 (a,x)\longrightarrow(b,y)
\]
is a morphism $f:a\to b$ satisfying $P(f)y\approx_a x$.
Its equality is the equality of the underlying arrows of $C$.  Pairing and
projection of codes are primitive recursive.  Identity and composition are
inherited from $C$.

\begin{lemma}
The displayed data define an arithmetic presentation of a category.
\end{lemma}
\begin{proof}
For identity, $P(\id_a)x\approx_a x$ is a presheaf law.  For composition,
let $f:(a,x)\to(b,y)$ and $g:(b,y)\to(c,z)$.  Then
\[
 P(g\circ f)z\approx_a P(f)(P(g)z)
 \approx_a P(f)y\approx_a x.
\]
Thus $g\circ f$ has the required element condition.  Congruence of the action
shows that changing an arrow to an equal arrow preserves this condition.
The equivalence, congruence, associativity, and identity laws are inherited
from $C$ with the indicated typing hypotheses.
\end{proof}

\begin{definition}[Chosen universal element]
A chosen universal element of $P$ is a valid pair $(A,x)$ together with a
fixed primitive recursive function $r(a,z)$ such that, whenever $z\in P(a)$,
\[
 r(a,z):a\to A,\qquad P(r(a,z))x\approx_a z,
\]
and every $f:a\to A$ satisfying $P(f)x\approx_a z$ satisfies
$f\sim r(a,z)$.  This is a terminal object of $\int P$ equipped with a
primitive recursive mediator function.
\end{definition}

\begin{theorem}[Representability by a chosen universal element]
A chosen universal element $(A,x,r)$ determines an isomorphism of presented
presheaves $h_A\cong P$.  Conversely, a presented isomorphism
$\theta:h_A\cong P$ determines such a chosen universal element.  The
constructions and their laws are schematic theorems of $\ipra$.
\end{theorem}
\begin{proof}
For the forward direction, put $\theta=\widehat{x}$ and
$\rho_a(z)=r(a,z)$.  The definition gives
$\theta_a(\rho_a(z))\approx_a z$.  It also implies
$\rho_a(\theta_a(f))\sim f$, since $f$ itself satisfies the relevant
universal condition.

We must check that $\rho$ is an extensional natural transformation, rather
than merely a family of inverse values.  If $z\approx_a z'$, then
$P(r(a,z))x\approx_a z'$; uniqueness for $z'$ gives
$r(a,z)\sim r(a,z')$.  Thus $\rho_a$ is extensional.  For $k:b\to a$,
\[
 P(r(a,z)\circ k)x
 \approx_b P(k)(P(r(a,z))x)
 \approx_b P(k)z.
\]
The uniqueness condition at $(b,P(k)z)$ gives
$r(a,z)\circ k\sim r(b,P(k)z)$.  This is naturality of $\rho:P\Rightarrow h_A$.
Therefore $\theta$ and $\rho$ form the required presheaf isomorphism.

Conversely, let $\theta:h_A\Rightarrow P$ and $\rho:P\Rightarrow h_A$ be
presented inverse transformations.  Set $x=\theta_A(\id_A)$ and
$r(a,z)=\rho_a(z)$.  Yoneda gives
$\theta_a(f)\approx_a P(f)x$.  Hence
\[
 P(r(a,z))x\approx_a\theta_a(\rho_a(z))\approx_a z.
\]
If $P(f)x\approx_a z$, the same Yoneda identity gives
$\theta_a(f)\approx_a z$.  Applying the extensional map $\rho_a$ and the
inverse law yields $f\sim\rho_a(z)=r(a,z)$.  This proves uniqueness.
\end{proof}

The adjective ``chosen'' is substantive.  A bare assertion that every object
of $\int P$ has some unique arrow to $(A,x)$ does not, by itself, supply a
single primitive recursive mediator function.  Ordinary intuitionistic logic
allows the corresponding local existential reasoning, but our claim that a
\emph{presented} inverse transformation has been constructed requires the
additional function $r$.  No choice or witness-extraction theorem is being
silently invoked.

\begin{theorem}[Uniqueness of representations, with the canonical arrow]
Suppose $\theta:h_A\cong P$ and $\psi:h_B\cong P$ are presented
isomorphisms.  There is a unique arrow $u:A\to B$, modulo hom equality,
such that
\[
 \psi_a(u\circ f)\approx_a\theta_a(f)
 \qquad(f:a\to A).
\]
It is an isomorphism, and its two directions are explicitly
\[
 u=\psi_A^{-1}(\theta_A(\id_A)),\qquad
 v=\theta_B^{-1}(\psi_B(\id_B)).
\]
\end{theorem}
\begin{proof}
The composite transformation $\psi^{-1}\theta:h_A\Rightarrow h_B$ is
induced by its evaluation at $\id_A$, which is the displayed $u$.
Therefore $\psi^{-1}\theta=h_u$ pointwise.  Composing pointwise with $\psi$
gives the required compatibility.  Any other compatible $u'$ satisfies
$h_{u'}=\psi^{-1}\theta=h_u$, so Yoneda faithfulness gives $u'\sim u$.
The analogous construction gives $h_v=\theta^{-1}\psi$.  These two
transformations are inverse; the preceding corollary gives
$v\circ u\sim\id_A$ and $u\circ v\sim\id_B$.
\end{proof}

\subsection{The covariant form}

For completeness, suppose $Q$ is a presented \emph{covariant} setoid-valued
functor: $f:a\to b$ gives $Q(f):Q(a)\to Q(b)$, with the covariant functor
laws.  Put $k_A(a)=\Hom_C(A,a)$.  The same statement holds with
$k_A$ in place of $h_A$:
\[
 \alpha\longmapsto\alpha_A(\id_A),\qquad
 x\longmapsto\bigl(f\longmapsto Q(f)x\bigr).
\]
To verify the nontrivial inverse directly, naturality at $f:A\to a$ gives
$Q(f)(\alpha_A(\id_A))\approx_a\alpha_a(f\circ\id_A)
\approx_a\alpha_a(f)$.  Naturality of the reconstruction is the covariant
composition law.  A universal element in this variance is a chosen
\emph{initial} object in the covariant category of elements, whose arrow
condition is $Q(f)x\approx_b y$.  Thus both variances are covered by explicit
formulas rather than an appeal to an unconstructed opposite functor category.

\section{Density and co-Yoneda by explicit normalization}
\label{sec:density}

There is a useful strengthening of the preceding analysis.  The density
formula for a presheaf can be proved in a form that avoids assuming the
existence of arbitrary colimits.  The key is that each apparent quotient
class has a canonical representative computable by one presheaf action.

Fix $b\in C$.  A generator over $b$ is a triple
\[
 (a,x,f),\qquad x\in P(a),\quad f:b\to a.
\]
Define its value by $\nu_b(a,x,f)=P(f)x\in P(b)$.  Put
\begin{equation}
 (a,x,f)\equiv_b(a',x',f')
 \quad\Longleftrightarrow\quad
 P(f)x\approx_b P(f')x'.
 \label{eq:density-equality}
\end{equation}
This is an arithmetic equivalence relation; it is not implemented by a
search for an equivalence class.  Let $D_P(b)$ denote this presented setoid.
An arrow $k:c\to b$ acts by
\[
 D_P(k)(a,x,f)=(a,x,f\circ k).
\]
The presheaf action law shows that this action respects
\eqref{eq:density-equality}, and the category laws prove its identity and
composition laws.  Thus $D_P$ is a presented presheaf.

\begin{theorem}[Explicit co-Yoneda normal form]
\label{thm:co-yoneda-normal}
There is a presented natural isomorphism $D_P\cong P$ with maps
\[
 \nu_b(a,x,f)=P(f)x,\qquad
 \iota_b(z)=(b,z,\id_b).
\]
Moreover, the relation \eqref{eq:density-equality} is exactly the least
setoid equivalence compatible with fibre and hom equalities and the elementary
identifications
\begin{equation}
 (a,x,f)\equiv_b(a',x',u\circ f)
 \quad\text{when }u:(a,x)\to(a',x')\text{ in }\int P.
 \label{eq:density-elementary}
\end{equation}
Here ``least'' is a schematic implication about any fixed arithmetic
relation satisfying these closure conditions; it does not quantify over a
set of all relations.
\end{theorem}
\begin{proof}
The maps have the required types.  The composite $\nu_b\iota_b$ sends $z$
to $P(\id_b)z\approx_b z$.  The other composite sends $(a,x,f)$ to
$(b,P(f)x,\id_b)$, whose value is
$P(\id_b)(P(f)x)\approx_b P(f)x$.  The definition of $\equiv_b$ therefore
makes this composite the identity modulo the setoid relation.

Naturality of $\nu$ is the identity
$P(f\circ k)x\approx_c P(k)(P(f)x)$.  For $\iota$, the images
$D_P(k)(b,z,\id_b)=(b,z,\id_b\circ k)$ and
$(c,P(k)z,\id_c)$ have equal values under $\nu_c$ by the identity and
composition laws; hence they are $\equiv_c$-equal.  This proves naturality.

We next identify the quotient relation intrinsically.  Fibre equality and
arrow equality plainly preserve values.  If $u:(a,x)\to(a',x')$, its
defining condition says $P(u)x'\approx_a x$.  Consequently
\[
 P(u\circ f)x'\approx_b P(f)(P(u)x')\approx_b P(f)x.
\]
Thus every elementary identification preserves $\equiv_b$.

Conversely, let $R$ be any equivalence relation on generators that respects
the stated fibre equalities and contains the elementary identifications.
For each generator $t=(a,x,f)$, the arrow
\[
 f:(b,P(f)x)\longrightarrow(a,x)
\]
is an arrow of $\int P$ by reflexivity.  Applying the elementary relation
to this arrow and $\id_b$ gives
\[
 (b,P(f)x,\id_b)\ R\ (a,x,f\circ\id_b)\ R\ (a,x,f),
\]
where the last step uses the specified compatibility with hom equality.
If $t\equiv_b t'$, their two values are fibre-equal, so the corresponding
canonical triples are $R$-equal.  Transitivity and symmetry now give $t R t'$.
This uses only a fixed finite sequence of equalities, regardless of the
numerical codes.  In particular, there is no appeal to induction on an
arbitrarily long chain of quotient generators.
\end{proof}

The customary coend notation for this result is
\[
 P(b)\cong\int^{a\in C}P(a)\times\Hom_C(b,a).
\]
In the present foundation, this expression means precisely the coded
setoid and normalization maps above.  It should not be read as a theorem
asserting all coends exist in a global category of sets.

\begin{theorem}[Schematic density as a universal cocone]
\label{thm:density-cocone}
For $(a,x)\in\int P$, let $j_{a,x}:h_a\Rightarrow P$ be the Yoneda
transformation $j_{a,x,b}(f)=P(f)x$.  This family is a cocone over the diagram
$(a,x)\mapsto h_a$, with an arrow $u$ acting by postcomposition $h_u$.
It has the following universal property.  For every presented presheaf $Q$
and every \emph{uniformly primitive recursive} family of transformations
$\lambda_{a,x}:h_a\Rightarrow Q$ satisfying
\begin{equation}
 \lambda_{a,x,b}(f)\approx_b
 \lambda_{a',x',b}(u\circ f)
 \quad\bigl(u:(a,x)\to(a',x')\bigr),
 \label{eq:density-cocone}
\end{equation}
there is a unique presented transformation $\delta:P\Rightarrow Q$ such
that $\delta j_{a,x}=\lambda_{a,x}$ pointwise.  It is explicitly
\begin{equation}
 \delta_b(z)=\lambda_{b,z,b}(\id_b).
 \label{eq:density-mediator}
\end{equation}
\end{theorem}
\begin{proof}
The cocone law for $j$ says
$P(f)x\approx_b P(u\circ f)x'$, and this was proved in the preceding theorem.
The displayed formula for $\delta$ is primitive recursive because the family
$\lambda$ is assumed uniform in all of its numerical arguments.  Its value
belongs to $Q(b)$ by typing of $\lambda_{b,z}$.

To check extensionality, suppose $z\approx_b z'$.  The identity arrow
$\id_b$ is an arrow $(b,z)\to(b,z')$ in $\int P$, since
$P(\id_b)z'\approx_b z$.  Applying the cocone equation at $\id_b$ gives
\[
 \lambda_{b,z,b}(\id_b)\approx_b
 \lambda_{b,z',b}(\id_b\circ\id_b)
 \approx_b\lambda_{b,z',b}(\id_b).
\]
The last equality uses the identity law and extensionality of
$\lambda_{b,z',b}$.  This is $\delta_b(z)\approx_b\delta_b(z')$.

For naturality, let $k:c\to b$.  The arrow
$k:(c,P(k)z)\to(b,z)$ lies in $\int P$.  The cocone condition gives
\[
 \delta_c(P(k)z)=\lambda_{c,P(k)z,c}(\id_c)
 \approx_c\lambda_{b,z,c}(k).
\]
Naturality of the fixed transformation $\lambda_{b,z}:h_b\Rightarrow Q$
at $k$, evaluated at $\id_b$, gives
\[
 \lambda_{b,z,c}(k)\approx_c
 Q(k)(\lambda_{b,z,b}(\id_b))=Q(k)(\delta_b(z)).
\]
These two equalities prove naturality of $\delta$.

To prove factorization, take $f:b\to a$ and $x\in P(a)$.  The arrow
$f:(b,P(f)x)\to(a,x)$ belongs to $\int P$, so the cocone condition gives
\[
 \delta_b(j_{a,x,b}(f))
 =\lambda_{b,P(f)x,b}(\id_b)
 \approx_b\lambda_{a,x,b}(f).
\]
Finally, let $\delta'$ be another factorization.  Applying its factorization
equation to the canonical generator $(b,z,\id_b)$ yields
\[
 \delta'_b(P(\id_b)z)\approx_b\lambda_{b,z,b}(\id_b)=\delta_b(z).
\]
The presheaf identity law and extensionality of $\delta'_b$ replace the
left side by $\delta'_b(z)$.  Thus $\delta'$ equals $\delta$ pointwise.
\end{proof}

\paragraph{Foundational content.}
The density theorem has an unbounded object predicate, but its proof still
requires no unbounded search, induction on that predicate, choice of
representatives, or generic recursion over formulas.  It constructs the
mediator by a single substitution into the \emph{supplied uniform} cocone
program.  Removing that uniformity assumption would remove the reason that
\eqref{eq:density-mediator} is primitive recursive.  This distinction is
precisely what allows a substantial part of presheaf theory to survive in
$\ipra$ while avoiding a claim that all classical presheaf categories have
been internalized.
